% !TeX program = xelatex
\documentclass[11pt,oneside]{article}
\usepackage{ProblemsetStyle}

\hypersetup{
    pdftitle={20260914 Problem set - Problems},
    pdfauthor={김정우},
    pdfsubject={Science Quiz mathematics practice problem set},
    pdfcreator={XeLaTeX}
}

\begin{document}

\problemsettitle{20260914 Problem set}

\begin{problem}{1}
집합
\[
    \{1,2,\ldots,10\}
\]
의 $3$-element subset을 uniformly random하게 하나 고른다. 선택된 세 정수 가운데 서로 연속한 두 정수가 존재하지 않을 확률을 구하시오.
\end{problem}

% Verified against OpenAI's official GPT-5.5 announcement dated April 23, 2026,
% and the accompanying paper "On the Ratio of R(k, l) and R(k, l+1)".
\begin{problem}{2}
2026년 4월 OpenAI는 internal GPT-5.5 system이 fixed integer $k\geq2$에 대하여
\[
    \lim_{\ell\to\infty}
    \frac{R(k,\ell+1)}{R(k,\ell)}
    =1
\]
이라는 longstanding asymptotic fact의 새로운 proof를 찾았고, 이후 Lean으로 검증되었다고 발표했다. 이 결과가 다루는 Ramsey numbers의 종류를 쓰시오.
\end{problem}

\begin{problem}{3}
다음 급수의 값을 구하시오.
\[
    \sum_{n=1}^{\infty}
    \arctan\!\left(\frac{1}{n^2+n+1}\right).
\]
\end{problem}

\begin{problem}{4}
다음 binomial coefficient를 나누는 가장 큰 $2$의 거듭제곱을 $2^m$이라 하자. $m$을 구하시오.
\[
    \binom{2026}{1013}.
\]
\end{problem}

% Verified against the official 2025 Breakthrough Prize in Mathematics citation.
\begin{problem}{5}
2025 Breakthrough Prize in Mathematics는 geometric Langlands program에 대한 foundational work, 특히 characteristic $0$에서 geometric Langlands conjecture의 proof에 중심적인 역할을 한 한 수학자에게 수여되었다. 수상자의 이름을 쓰시오.
\end{problem}

\begin{problem}{6}
세 변의 길이가 $13$, $14$, $15$인 삼각형의 circumradius를 $R$, inradius를 $r$이라 하자. 다음 값을 구하시오.
\[
    R-2r.
\]
\end{problem}

\newpage

\begin{problem}{7}
$4\times4$ real matrix $\mathbf{A}$가 nilpotent이고
\[
    \operatorname{rank}(\mathbf{A})=3
\]
이라고 하자. $\operatorname{rank}(\mathbf{A}^2)$을 구하시오.
\end{problem}

% Verified against the International Mathematical Union's July 2026 memorial article.
\begin{problem}{8}
2026년 3월 18일 $94$세로 별세한 1970 Fields Medalist의 이름을 쓰시오. 이 수학자는 characteristic $0$에서 algebraic varieties의 resolution of singularities를 arbitrary dimension에 대해 증명한 업적으로 잘 알려져 있다.
\end{problem}

\begin{problem}{9}
다음 값을 구하시오.
\[
    \sum_{k=1}^{100}\left\lfloor\sqrt{k}\right\rfloor.
\]
\end{problem}

\begin{problem}{10}
complete graph $K_6$에서 지정된 edge 하나를 제거하여 얻은 graph의 spanning tree 개수를 구하시오.
\end{problem}

\end{document}
